\documentclass[11pt]{article}
\usepackage[margin=1in]{geometry}
\usepackage{amsmath,amssymb}
\usepackage{booktabs}
\usepackage{hyperref}

\title{Scenario 1 --- Shrinking-Horizon MPC\\
       for Mobile Charging Station Dispatch\\[2pt]
       \large Formal Model (Approach 2 tree)}
\author{Sections 1--11: the deterministic MILP (Approach 0's baseline, and the per-scenario
building block of Approach 2's stochastic extension). Section~\ref{sec:stochastic}: the
stochastic (scenario-based) extension itself --- Approach 2's actual closed loop.}
\date{}

\newcommand{\dt}{\Delta t}

\begin{document}
\maketitle

Formal companion to the code in \texttt{Shrinking\_Horizon/code/} (the window MILP is
built in \texttt{3\_MCSModel.jl}; the closed loop in \texttt{4\_MPCLoop.jl}). Plain-English
descriptions and the data schema are in \texttt{README.md}; a code audit is in
\texttt{constraints\_code\_vs\_model.txt}. Sections~\ref{sec:sets}--\ref{sec:fixedH} below
document the \textbf{deterministic} model, \texttt{build\_window\_model}: solved \emph{once},
then replayed open-loop, by Approach~0 (the one-shot baseline used to measure the value of
re-planning), and reused --- unchanged, per scenario --- as the building block of Approach~2's
own stochastic model in Section~\ref{sec:stochastic}, which is what this tree's actual closed
loop (\texttt{run\_mpc}) solves every interval. Approach~0 is replayed under \emph{two} plant
modes --- a deterministic one, in which the realized powers equal the planning mean so the
outcome is the MILP's own optimum, and a stochastic one --- so that plan drift and re-planning
value can be read apart. Plant modes are defined in \S\ref{sec:plant}.

The controller is a \textbf{single-day, full-24\,h shrinking horizon}. Each 15-minute step
re-solves a MILP over the \emph{entire remaining day} $[k_0,\dots,n_{\text{day}}]$; only the
first interval is applied, and the window shrinks as the day advances. There is one
optimisation per step and one horizon: the overnight MCS recharge is scheduled \emph{inside}
this MILP (there is no separate phase). The activity powers are a \textbf{fixed} calibrated
estimate; \textbf{Approach 0} plans on the mean directly, while \textbf{Approach 2} samples a
scenario set from this same estimate every re-solve (Section~\ref{sec:stochastic}). Either way,
the stochastic plant draws realized powers from a
pool pre-sampled from $\mathcal{N}(\mu,\sigma)$, one draw per (excavator, activity) occurrence
(shared across every approach that simulates the plant, so their realizations are comparable).

%======================================================================
\section{Sets and indices}
\label{sec:sets}
\begin{tabular}{ll}
$m \in \mathcal{M}$ & mobile charging stations (MCS) \\
$e \in \mathcal{E}$ & construction EVs (excavators) \\
$i,j \in \mathcal{N}$ & nodes; $\mathcal{N}_g$ grid, $\mathcal{N}_c$ sites \\
$a \in \mathcal{B}$ & activities $=\{\text{dig},\text{load},\text{trv},\text{idle}\}$ \\
$k$ & interval $1,\dots,n_{\text{day}}$; window $[k_0,n_{\text{day}}]$ \\
$\mathcal{T}_b$ & interval boundaries of the window \\
\end{tabular}

\medskip
$n_{\text{day}} = n_{\text{int}} = 96$ is the \textbf{full 24\,h} day at $\dt=0.25$\,h
(08:00 $\to$ 08:00 next day). $A_{i,e}=1$ iff CEV $e$ works at site $i$. Work availability
$R_{i,e,k}$ is $0$ outside the shift, but the MCS/CEVs may charge and the MCS may travel at
any interval.

%======================================================================
\section{Parameters (selected)}
\begin{tabular}{ll}
$\dt$ & interval length (h), $0.25$ \\
$\lambda^{\text{buy}}_{k},\lambda^{\text{CO2}}_{k}$ & price (\$/kWh) / carbon intensity \\
$\overline{SOE}_\bullet,\underline{SOE}_\bullet,SOE^{\text{ini}}_\bullet$ & energy cap / floor / start (MCS, CEV) \\
$CH_m,DCH_m,DCH^{\text{plug}}_m,C^{\text{plug}}_m,\eta_m,CH^{\text{CEV}}_e$ & MCS / CEV charge-discharge parameters \\
$p_a$ & activity power (kW); fixed calibrated mean, MILP plans on it \\
$H^{\text{dig}}_i,H^{\text{load}}_i$ & lumpsum dig / load requirement (hours) at site $i$ \\
$\tau_{i,j},k_{\text{trv}},R_{i,e,k}$ & travel time / energy / work availability \\
$\rho_{\text{miss}},\rho_{\text{lab}},\lambda^{\text{NC}},\lambda^{\text{OP}}$ & penalties \\
$\text{scale},w_{\text{pt}}=4,r$ & precedence ratio / travels-per-work / rest cap \\
\end{tabular}

%======================================================================
\section{Decision variables}
\textbf{Continuous ($\ge0$):} $P^{\text{ch}}_{m,i,k}$, $P^{\text{dch}}_{m,i,k}$,
$P^{\text{MCS-CEV}}_{m,i,e,k}$, $P^{\text{work}}_{i,e,k}$, totals
$P^{\text{ch,tot}}_{m,k},P^{\text{dch,tot}}_{m,k}$, travel energy $L_{m,i,j,k}$, state
$SOE^{\text{MCS}}_{m,t},SOE^{\text{CEV}}_{e,t}$, peaks $P^{\text{NC}},P^{\text{OP}}$, and the
single work-shortfall slack $s^{\text{miss}}_{i,a}$.

\smallskip
\textbf{Binary:} $u_{e,i,a,k}$, $\mu_{i,e,k}$, $\rho_{m,i,e,k}$, $z_{m,i,k}$,
$g^{\text{ch}}_{m,i,k}$, $x_{m,i,j,k}$, $y_{m,i,j,k}$, $\beta^{\text{arr}},\beta^{\text{dep}}$.

%======================================================================
\section{Objective (Eq.\ 1)}
\begin{align}
\min\ &\sum_{m,k}\lambda^{\text{buy}}_{k}P^{\text{ch,tot}}_{m,k}\dt
      +\sum_{m,k}\tfrac{c_{\text{ton}}}{1000}\lambda^{\text{CO2}}_{k}P^{\text{ch,tot}}_{m,k}\dt
      +\rho_{\text{miss}}\!\!\sum_{i\in\mathcal{N}_c,a}\!\! s^{\text{miss}}_{i,a}\notag\\
     &+\lambda^{\text{NC}}P^{\text{NC}}+\lambda^{\text{OP}}P^{\text{OP}}
      +\rho_{\text{lab}}\dt\!\!\sum_{m,i,j,k}\!\! y_{m,i,j,k}.
\end{align}

%======================================================================
\section{Constraints}

\paragraph{Power aggregation \& flow (HARD).}
$P^{\text{ch,tot}}_{m,k}=\sum_{i\in\mathcal{N}_g}P^{\text{ch}}_{m,i,k}$,
$P^{\text{dch,tot}}_{m,k}=\sum_{i\in\mathcal{N}_c}P^{\text{dch}}_{m,i,k}$;
$P^{\text{dch}}_{m,i,k}=\sum_e P^{\text{MCS-CEV}}_{m,i,e,k}\le DCH_m z_{m,i,k}$;
$P^{\text{MCS-CEV}}_{m,i,e,k}\le DCH^{\text{plug}}_m\rho_{m,i,e,k}$;
$\sum_m P^{\text{MCS-CEV}}_{m,i,e,k}\le CH^{\text{CEV}}_e\mu_{i,e,k}$.
Grid exclusivity: $P^{\text{ch}}_{m,i,k}\le CH_m g^{\text{ch}}_{m,i,k}$,
$g^{\text{ch}}\le z$, $\sum_m g^{\text{ch}}_{m,i,k}\le 1$.

\paragraph{State of energy (HARD).}
\begin{align}
SOE^{\text{MCS}}_{m,k+1}&=SOE^{\text{MCS}}_{m,k}+\eta_m P^{\text{ch,tot}}_{m,k}\dt-\tfrac{1}{\eta_m}P^{\text{dch,tot}}_{m,k}\dt-L^{\text{tot}}_{m,k},\\
SOE^{\text{CEV}}_{e,k+1}&=SOE^{\text{CEV}}_{e,k}+\textstyle\sum_{m,i}P^{\text{MCS-CEV}}_{m,i,e,k}\dt-\sum_i P^{\text{work}}_{i,e,k}\dt,\\
\underline{SOE}_\bullet&\le SOE_\bullet\le\overline{SOE}_\bullet .
\end{align}

\paragraph{Terminal energy targets (Eq.\ 8a / 8b, HARD).}
Applied at the final window boundary \emph{when that boundary is the end of the day}, i.e.
only when $\max K=n_{\text{day}}$. On the shrinking horizon this holds for every window, so
the rules always bind; under the alternative fixed-$H$ lookahead they are absent from all
but the last $H$ windows and nothing replaces them (see \S\ref{sec:fixedH}).
\begin{align}
SOE^{\text{MCS}}_{m,n_{\text{day}}+1}&= SOE^{\text{MCS,ini}}_m && \text{(8a, exact; overnight refill inside this MILP)},\\
SOE^{\text{CEV}}_{e,n_{\text{day}}+1}&\ge SOE^{\text{CEV,ini}}_e && \text{(8b, floor; overcharging allowed)}.
\end{align}
The CEV target is a \emph{floor} rather than Avik's exact equality: a CEV cannot discharge, so
under the stochastic plant an equality can become unrecoverable; the floor keeps the terminal
reachable while guaranteeing the fleet ends at least as charged as it began.

\paragraph{Plugging, presence \& routing (Eq.\ 10, HARD).}
$\rho\le A$, $\rho\le z$, $\sum_e \rho_{m,i,e,k}\le C^{\text{plug}}_m$. Presence partition
$\sum_i z_{m,i,k}+\sum_{i\ne j}y_{m,i,j,k}=1$; $\beta^{\text{dep}}_{m,i,k}=\sum_j x_{m,i,j,k}$;
$\beta^{\text{arr}}$ from finishing trips (or a carried-in arrival);
$\beta^{\text{arr}}-\beta^{\text{dep}}=z_k-z_{k-1}$; $\beta^{\text{arr}}+\beta^{\text{dep}}\le1$;
flow balance (generalised for the carried-in start). Trip on $i\!\to\!j$ lasts $\tau_{i,j}$ and
costs $L=k_{\text{trv}}\dt\,y$. \textbf{(10e)} MCS parked at a grid node by day-end:
$\sum_{i\in\mathcal{N}_g}z_{m,i,n_{\text{day}}}=1$.

\paragraph{Activity scheduling (Eq.\ 11, HARD).}
$\sum_a u_{e,i,a,k}=A_{i,e}$; $u\le A$; $\mu_{i,e,k}\le u_{e,i,\text{idle},k}$;
$P^{\text{work}}_{i,e,k}\le R_{i,e,k}A_{i,e}(1-\mu_{i,e,k})$;
$P^{\text{work}}_{i,e,k}=\sum_a p_a u_{e,i,a,k}$ (Eq.\ 5e), with $p_{\text{idle}}=0$.
(The code uses this 4-activity encoding; Avik's equivalent 3-activity form is
$\sum_a u\le1$, $\sum_a u+\mu\le1$.)

\paragraph{Work quota (Eq.\ 12c, SOFT shortfall + implicit HARD cap) --- lumpsum.}
Written in the code as a balance \emph{equality} against the work still outstanding
($H^{\bullet}_i$ below is the \emph{remaining} requirement $\text{rem}^{\bullet}_i$ at the
start of the window, not the whole-day figure):
\begin{equation}
\dt\!\!\sum_{e,k}\!u_{e,i,\text{dig},k}+s^{\text{miss}}_{i,\text{dig}}=H^{\text{dig}}_i,\qquad
\dt\!\!\sum_{e,k}\!u_{e,i,\text{load},k}+s^{\text{miss}}_{i,\text{load}}=H^{\text{load}}_i,\qquad
s^{\text{miss}}\ge0 .
\end{equation}
The shortfall $s^{\text{miss}}$ is soft (priced at $\rho_{\text{miss}}$ in Eq.~1), but because
$s^{\text{miss}}\ge0$ the equality also implies the \textbf{hard} upper bound
$\dt\sum_{e,k}u_{e,i,a,k}\le H^{a}_i$: a CEV cannot do more dig/load work than the requirement
still outstanding. This is the shrinking-horizon analogue of an explicit no-working-ahead rule,
obtained implicitly, and it is re-imposed against the \emph{remaining} requirement at every
re-solve. (Earlier revisions of this document wrote the constraint as
$s^{\text{miss}}\ge H-\dt\sum u$, which permits overshoot; that did not match the code.)

\paragraph{Precedence (Eq.\ 12d, HARD).}
Raw interval counts, seeded by realized work carried in:
$\sum_{e,\tau\le k}u_{e,i,\text{load},\tau}\le\text{scale}\sum_{e,\tau\le k}u_{e,i,\text{dig},\tau}$.

\paragraph{Rest rule (Eq.\ 12e, HARD).}
With $r=\mathrm{round}(t_{\text{rest}}/\dt)$ (=4 at the defaults $t_{\text{rest}}=1$\,h,
$\dt=0.25$\,h; the code rounds, it does not take a ceiling),
$\sum_{k'=k}^{k+r}\sum_{a\in\{\text{dig,load,trv}\}}u_{e,i,a,k'}\le r$, \emph{seeded from the
applied Work/Break history} so a work-run cannot leak across the every-15-min re-solves.

\paragraph{Travel pacing (Eq.\ 13, HARD, no tolerance).}
With $w_{\text{pt}}=4$, for each $(i,e)$ and $k$, on cumulative travel $V$ vs cumulative
useful work $W$ (dig+load):
\begin{equation}
W(k)-w_{\text{pt}}\ \le\ w_{\text{pt}}\,V(k)\ \le\ W(k),
\end{equation}
where $V$ and $W$ are raw \emph{applied interval counts} (whether the $u$ indicator fired that
interval), read off \texttt{cum\_trv\_cnt\_e}/\texttt{cum\_work\_cnt\_e}, not hours. A battery-
shortage-capped interval (e.g.\ 14 of 15 minutes realized) still counts as one full travel/work
interval, so the sub-interval rounding residue that used to strand the band on an integer boundary
(see \texttt{cum\_dig\_e}/\texttt{cum\_load\_e} in the precedence rule) cannot occur here, and no
tolerance $\tau$ is needed. Precedence (above) still uses realized hours; only pacing switched to
interval counts. Applies to both the deterministic and scenario-linked stochastic builder.

\paragraph{Demand peaks (E1, HARD).}
$P^{\text{NC}}\ge$ carried-in peak and $\ge\sum_m P^{\text{ch,tot}}_{m,k}$ for all $k$;
$P^{\text{OP}}$ likewise on on-peak $k$ only.

%======================================================================
\section{Single horizon, fixed model}
There is no overnight ``Phase~2'': the terminal rules (8a/8b) close the MCS energy cycle inside
the one 24\,h MILP. The activity powers $p_a=\mu_a$ are the posterior \emph{mean} of a Bayesian
fit performed \textbf{offline, once, before the loop} by step~0 (\texttt{0\_Regression.jl}): it
reads the raw task-recording data, builds cumulative-$\Delta$SoC energy-balance equations
$A x = b$, and fits $b\sim\mathcal{N}(Ax,s)$ with $\mathrm{TruncatedNormal}(\mu_0,\sigma_0;\ge0)$
power priors and a half-normal noise (pooled NUTS chains). It writes both the mean $\mu$ and a
\textbf{per-activity} standard deviation $\sigma_a$ to \texttt{parameters.csv}. The MILP plans on
$\mu$; the closed-loop plant draws each realized power from a pool pre-sampled from
$\mathcal{N}(\mu_a,\sigma_a)$ (truncated at $0$, idle pinned to $0$) once, before the run, and
consumes the next unused sample for a $(\text{CEV},a)$ pair each time that pair actually occurs
--- not a fresh draw every interval. The model is \emph{not} re-fit online during the day.

%======================================================================
\section{Plant modes: what the disturbance actually is}
\label{sec:plant}
The MILP above is identical in every mode; only what the \emph{plant} does with the applied
plan differs. Both \texttt{run\_mpc} (Approach~1) and \texttt{run\_one\_shot} (Approach~0)
take a \texttt{plant} switch:

\medskip
\begin{tabular}{ll}
\texttt{:sampled} & realized power $\hat p_{e,a}$ = the next unused draw from the pre-sampled\\
                  & pool; the within-interval activity split may be randomised. Realized\\
                  & $\ne$ planned, so the day drifts. \emph{(default)}\\[4pt]
\texttt{:mean}    & realized power $\hat p_{e,a}=\mu_a$, the same mean the MILP planned on, and\\
                  & each interval realizes its single planned activity in full. Realized\\
                  & $=$ planned \emph{exactly}. No pool sample is consumed and no cursor\\
                  & advances, so a \texttt{:mean} run cannot perturb a \texttt{:sampled} run\\
                  & sharing the same pool, and needs no seed to be reproducible.\\
\end{tabular}

\medskip
Approach~0 under \texttt{:mean} is therefore the whole-day MILP's \emph{own optimum} evaluated
over the full day --- the deterministic reference, with nothing stochastic anywhere --- and
the difference between the modes separates two effects that a single headline figure conflates:
\begin{align*}
\text{A0}(\texttt{:mean})\to\text{A0}(\texttt{:sampled})\ &=\ \text{cost of plan drift with no feedback},\\
\text{A0}(\texttt{:sampled})\to\text{A1}\ &=\ \text{value of re-planning against that drift}.
\end{align*}

\paragraph{Scope of the disturbance.} Only the CEV side is stochastic. The MCS state is
advanced by the \emph{planned} value $SOE^{\text{MCS}}_{m,k_0+1}$ from the solved window, so
grid draw, travel loss and hence every cost term follow the plan exactly; the disturbance
enters through CEV depletion and reaches the objective only indirectly, through what the next
re-solve decides to do about it.

\paragraph{State guard.} The realized dig/load/travel duration for each CEV is capped by the
energy actually available before hitting $\underline{SOE}_e$: if the sampled draw would push the
balance below the floor, the executed duration is scaled down to exactly what the floor allows,
crediting only the affordable fraction of that interval's work (the remainder becomes idle) before
the SOE update is applied. This is a physical correction, not a numerical guard — no energy is
created or destroyed, and \texttt{rem\_dig}/\texttt{rem\_load} are only reduced by what was truly
completed. A residual \texttt{clamp} to $[\underline{SOE}_e,\overline{SOE}_e]$ remains afterward as
a pure safety net and should not bind under normal operation. Capping events are counted and
reported per run ($\texttt{n\_capped}$), because they are invisible to the MILP's own feasibility
status.

%======================================================================
\section{The fixed-\texorpdfstring{$H$}{H} lookahead is not this controller}
\label{sec:fixedH}
The code exposes a switch that replaces the shrinking window $[k_0,n_{\text{day}}]$ with a fixed
lookahead $[k_0,\min(k_0+H-1,n_{\text{day}})]$. The terminal rules (8a), (8b) and (10e) are conditioned on the window actually reaching
day-end, so under a fixed $H$ they are absent from all but the final $H$ windows, and
\textbf{nothing replaces them} --- there is no terminal cost and no value-function approximation
standing in for the discarded tail. The resulting controller is myopic and will discharge the MCS
early with no obligation to restore it. The fixed-$H$ mode is provided for experimentation only;
a terminal cost must be added before it is used for reported results.

%======================================================================
\section{Approach 2: the stochastic (scenario-based) extension}
\label{sec:stochastic}
Everything above is the deterministic window MILP exactly as Approach~1 (and Approach~0's
one-shot baseline) solve it, and \texttt{build\_window\_model} in \texttt{3\_MCSModel.jl}
implements it unchanged. Approach~2's controller solves a second, sibling model,
\texttt{build\_window\_model\_stochastic}, in its place. This section states the formal
difference only; see \texttt{Understanding\_Stochastic\_MPC.md} for the conceptual
development and a worked numerical example, and \S13 of \texttt{README.md} for the
implementation-level detail (RNG streams, the scenario-1 read-back view, etc.).

\paragraph{Scenario set.} At each re-solve, instead of the single mean vector $\mu$, the
controller samples $S$ power vectors $\{\mu^{(s)}\}_{s=1}^{S}$ from the same fitted posterior,
$\mu^{(s)}_a \sim \mathcal{N}(\mu_a,\sigma_a)^+$ (clipped at 0), independently per scenario,
via a dedicated RNG stream disjoint from the plant's own draws. Each scenario carries equal
weight $\pi_s = 1/S$ ($S=5$ by default). This is the discrete empirical approximation of the
posterior used throughout \texttt{Understanding\_Stochastic\_MPC.md}'s Level 4/5 treatment.

\paragraph{Scenario-indexed model.} Every decision variable $X_{\dots,k}$ defined above
gains a scenario index, $X_{\dots,k}^{(s)}$, $s=1,\dots,S$, and every constraint of
the sections above is imposed \emph{for every $s$ individually} — HARD in every
scenario, not on average. The objective becomes the probability-weighted average
\[
  \min \ \sum_{s=1}^{S} \pi_s \, J^{(s)}\bigl(x^{(s)}\bigr),
\]
where $J^{(s)}$ is exactly the Eq.~(1) objective evaluated against scenario $s$'s copy of the
variables, and the sole substitution inside each scenario's copy is $\mu \to \mu^{(s)}$
wherever the planned power feeds a work-power identity (Eq.~5e).

\paragraph{Non-anticipativity.} Let $k_0=\text{first}(K)$ be the window's first (``here-and-now'')
interval. For every control family $X\in\{u,\,\mu^{\text{plug}},\,\rho,\,z,\,g^{ch},\,x^{\text{trip}},\,
y^{trv},\,\beta^{arr},\,\beta^{dep},\,P^{ch}_{MCS},\,P^{MCS}_{CEV},\,P^{dch}_{MCS},\,P^{ch}_{tot},\,
P^{dch}_{tot},\,L^{trv},\,L^{trv}_{tot}\}$ (using $\mu^{\text{plug}}$ to disambiguate the
charge-permission binary from the posterior mean $\mu$):
\[
  X^{(s)}_{\dots,k_0} \;=\; X^{(1)}_{\dots,k_0} \qquad \forall\, s = 2,\dots,S.
\]
This is the discrete-scenario encoding of the filtration constraint ``$x_{k_0}$ must be
$\mathcal F_{k_0}$-measurable'' (Level 5 of \texttt{Understanding\_Stochastic\_MPC.md}): the
action taken this instant cannot depend on information (which scenario is real) that does not
yet exist. \textbf{Deliberately excluded} from this set: $P^{work}_{i,e,k_0}$ and
$SOE^{CEV}_{e,k_0+1}$. Since $P^{work}_{i,e,k_0}=\sum_a \mu^{(s)}_a\, u_{e,i,a,k_0}$ (Eq.~5e,
scenario-indexed) and $u$ \emph{is} tied while $\mu^{(s)}$ is not, the chosen activity is shared
but its realized energy cost is scenario-dependent by construction — the entire mechanism by
which a single first-stage decision carries an uncertain consequence. $SOE^{CEV}$ inherits this
at $k_0+1$ through its recursion (Eq.~7). $SOE^{MCS}$ is not explicitly tied but is implied
equal across scenarios at $k_0+1$, since every quantity in its recursion at $k_0$ ($P^{ch}_{tot}$,
$P^{dch}_{tot}$) is itself tied — consistent with the earlier remark that only the CEV side is
stochastic'' (Plant modes section). $s^{\text{miss}}$, $P^{peak}_{NC}$, $P^{peak}_{OP}$ are whole-window aggregates over
the (scenario-dependent) trajectory and are left free.

\paragraph{Recovered special case.} At $S=1$ this collapses exactly to the deterministic model
of the sections above with $\mu^{(1)}=\mu$: no non-anticipativity constraints are
generated (there is nothing to tie), and the objective is a single scenario's cost with weight 1.

\end{document}
